% Preâmbulo comum aos documentos do Grupo VI (classe abnTeX2, normas ABNT).
% Compilar com Tectonic (XeTeX): ver gerar_pdf.sh.
\documentclass[12pt,openright,oneside,a4paper,chapter=TITLE,english,brazil]{abntex2}

% ---------------------------------------------------------------- Fontes
\usepackage{fontspec}
\setmainfont{texgyretermes}[
  Extension=.otf,
  UprightFont=*-regular, BoldFont=*-bold,
  ItalicFont=*-italic, BoldItalicFont=*-bolditalic]
\setmonofont{DejaVuSansMono}[
  Extension=.ttf, BoldFont=*-Bold, Scale=0.82]
\setsansfont{texgyreheros}[
  Extension=.otf,
  UprightFont=*-regular, BoldFont=*-bold,
  ItalicFont=*-italic, BoldItalicFont=*-bolditalic]
\usepackage{amsmath}

% ---------------------------------------------------------------- Pacotes
\usepackage{indentfirst}
\usepackage{graphicx}
\usepackage{xcolor}
\usepackage{booktabs}
\usepackage{array}
\usepackage{tabularx}
\usepackage{longtable}
\usepackage{xltabular}
\usepackage{multirow}
\usepackage[section]{placeins}
\usepackage{enumitem}
\usepackage{listings}
\usepackage{tikz}
\usepackage{adjustbox}
\usetikzlibrary{arrows.meta,positioning,fit,calc,shapes.geometric,backgrounds}
\usepackage[alf,abnt-emphasize=bf]{abntex2cite}

\graphicspath{{../../backend/resultados/}}

% ---------------------------------------------------------------- Aparência
\definecolor{cinzacodigo}{gray}{0.96}
\definecolor{cinzaborda}{gray}{0.75}
\lstset{
  language=Python,
  basicstyle=\ttfamily\small,
  keywordstyle=\bfseries,
  commentstyle=\itshape\color{black!65},
  stringstyle=\color{black!80},
  backgroundcolor=\color{cinzacodigo},
  frame=single, rulecolor=\color{cinzaborda},
  framesep=4pt, xleftmargin=6pt, xrightmargin=6pt,
  columns=fullflexible, keepspaces=true,
  showstringspaces=false, breaklines=true,
  numbers=none, aboveskip=0.8em, belowskip=0.8em,
}
\lstdefinestyle{json}{language={}, basicstyle=\ttfamily\footnotesize, keywordstyle={}}
\renewcommand{\lstlistingname}{Código}
\AtBeginDocument{\renewcommand{\thelstlisting}{\arabic{lstlisting}}}
\renewcommand{\lstlistlistingname}{Lista de códigos}

\setlist{itemsep=0.2em, topsep=0.3em}
\setlist[itemize]{label=--}
\newcolumntype{L}[1]{>{\raggedright\arraybackslash}p{#1}}
\newcolumntype{Y}{>{\raggedright\arraybackslash}X}
\newcolumntype{C}[1]{>{\centering\arraybackslash}p{#1}}
\renewcommand{\arraystretch}{1.25}
\setlength{\arrayrulewidth}{0.5pt}
% Quadros (ABNT): texto, bordas fechadas. Tabelas (IBGE): dados numéricos, laterais abertas.
% Definição do ambiente quadro conforme a documentação do abnTeX2.
\newcommand{\quadroname}{Quadro}
\newcommand{\listofquadrosname}{Lista de quadros}
\newcommand{\quadroautorefname}{Quadro}
\newfloat[chapter]{quadro}{loq}{\quadroname}
\newlistof{listofquadros}{loq}{\listofquadrosname}
\newlistentry{quadro}{loq}{0}
\setfloatadjustment{quadro}{\centering}
\counterwithout{quadro}{chapter}
\renewcommand{\cftquadroname}{\quadroname\space}
\renewcommand*{\cftquadroaftersnum}{\hfill--\hfill}
\setfloatlocations{quadro}{hbtp}

% Fonte de figuras e tabelas produzidas pelo grupo
\newcommand{\fonteautores}{\fonte{Elaborado pelos autores (2026).}}
% Diagrama TikZ: \diagrama{arquivo} reduz a figura à largura do texto, se preciso
\newcommand{\diagrama}[1]{\begin{adjustbox}{max width=\textwidth}\input{figuras/#1}\end{adjustbox}\par}
% Nome de arquivo, função ou campo do código
% (permite quebra de linha depois de \_ e de /)
\newcommand{\codsublinhado}{\textunderscore\allowbreak}
\newcommand{\cod}[1]{{\ttfamily\let\_\codsublinhado#1}}

% ---------------------------------------------------------------- PDF
\makeatletter
\hypersetup{
  pdftitle={\@title},
  pdfauthor={Grupo VI -- Sistemas Distribuídos -- UNIFESSPA},
  pdfsubject={Agregador de Preços de Eletrônicos},
  pdfcreator={LaTeX com abnTeX2},
  colorlinks=false, hidelinks, bookmarksdepth=4}
\makeatother

\setlength{\parindent}{1.25cm}
\setlength{\parskip}{0.2cm}
\setlength{\emergencystretch}{3em}

% Títulos em Times, negrito (seções primárias em caixa alta, como pede a NBR 6024)
\renewcommand{\ABNTEXchapterfont}{\rmfamily\bfseries}
\renewcommand{\ABNTEXchapterfontsize}{\normalsize}
\renewcommand{\ABNTEXsectionfont}{\rmfamily\bfseries}
\renewcommand{\ABNTEXsectionfontsize}{\normalsize}
\renewcommand{\ABNTEXsubsectionfont}{\rmfamily\bfseries}
\renewcommand{\ABNTEXsubsectionfontsize}{\normalsize}
\renewcommand{\ABNTEXsubsubsectionfont}{\rmfamily\itshape}
\renewcommand{\ABNTEXsubsubsectionfontsize}{\normalsize}

% ---------------------------------------------------------------- Capa e folha de rosto
% Versões próprias (NBR 14724): instituição, autores, título, subtítulo, local e ano.
\newcommand{\subtitulodoc}{}
\newcommand{\definirsubtitulo}[1]{\renewcommand{\subtitulodoc}{#1}}
\makeatletter
\newcommand{\capagrupo}{%
  \begin{capa}
    \centering
    {\bfseries\imprimirinstituicao\par}
    \vfill
    {\bfseries\large\MakeUppercase{\imprimirtitulo}\par}
    \ifx\subtitulodoc\@empty\else
      \vspace{0.4cm}{\large\subtitulodoc\par}\fi
    \vfill
    {\imprimirlocal\par}
    {\imprimirdata\par}
    \vspace*{0.5cm}
  \end{capa}}
\newcommand{\folhaderostogrupo}{%
  \begin{folhaderosto}
    \centering
    \vspace*{\fill}
    {\bfseries\large\MakeUppercase{\imprimirtitulo}\par}
    \ifx\subtitulodoc\@empty\else
      \vspace{0.4cm}{\large\subtitulodoc\par}\fi
    \vspace*{\fill}
    \hspace{.45\textwidth}%
    \begin{minipage}{.5\textwidth}
      \SingleSpacing\small
      \imprimirpreambulo
      \par\vspace{0.6cm}
      % nomes na mesma linha, separados por vírgula (\\ de \autor vira ",")
      \textbf{Discentes:} {\def\\{,}\imprimirautor}.\par
      \vspace{0.6cm}
      \textbf{\imprimirorientadorRotulo}~\imprimirorientador
    \end{minipage}\par
    \vspace*{\fill}
    {\imprimirlocal\par}
    {\imprimirdata\par}
    \vspace*{0.5cm}
  \end{folhaderosto}}
\makeatother
\newcommand{\listafiguras}{%
  \pdfbookmark[0]{\listfigurename}{lof}%
  \listoffigures*%
  \cleardoublepage}
\newcommand{\listatabelas}{%
  \pdfbookmark[0]{\listtablename}{lot}%
  \listoftables*%
  \cleardoublepage}
% Página pré-textual com título centralizado e sem número de seção
\newcommand{\paginapretextual}[1]{%
  \pdfbookmark[0]{#1}{\detokenize{#1}}%
  \begin{center}{\ABNTEXchapterfont\MakeUppercase{#1}}\end{center}\vspace{0.8cm}}
\newcommand{\listaquadros}{%
  \pdfbookmark[0]{\listofquadrosname}{loq}%
  \listofquadros*%
  \cleardoublepage}
% Legenda de quadro fora de float (para quadros longos com xltabular)
\newfixedcaption{\legendaquadro}{quadro}
\newcommand{\sumario}{%
  \pdfbookmark[0]{\contentsname}{toc}%
  \tableofcontents*%
  \cleardoublepage}
